\scanpage{31}
% archive foliation: 30
(Si $A$ est réduite à un objet, \inserted{$A$ \unsure{seul}}, l'anneau
des endom. $k$ de celui-ci, on
trouve une équivalence avec la
cat. des $k$-modules à droite
\uline{libres de type fini}, par $X\longmapsto\operatorname{Hom}(A,X)$)
\par\bigskip\hrule\bigskip
\unsure{Transformant|Transformons} par dualité les énoncés.
Si $\forall$ objet de $C$ est $\varprojlim$
d'objets de $A$, alors le foncteur
canonique
\[
\rho:C^\circ\longrightarrow\widehat{A^\circ}
\quad(\text{resp. }\widehat{A^\circ}_{ab})
\]
\margin{$\underline{\operatorname{Hom}}(A,\mathrm{Ens})$
(resp. $\underline{\operatorname{Hom}}(A,\mathrm{Ab})$)}
vers les foncteurs
\uline{covariants} sur $A$ à valeurs dans
$\mathrm{Ens}$ (resp. $\mathrm{Ab}$) est pl. fidèle. Si
de plus, $\forall$ objet de $C$ est
$\varinjlim$ de certains types d'objets
de $C$, alors l'image essentielle est
contenue dans la sous-cat. pleine
de $\widehat{A^\circ}$ (resp. $\widehat{A^\circ}_{ab}$) formée des
foncteurs $A\longrightarrow\mathrm{Ens}$ (resp. $A\longrightarrow\mathrm{Ab}$)
qui sont $\varprojlim$ de ces types de
foncteurs représentables --- et
l'image essentielle est $\simeq$ égale à
cette dernière catégorie si
